\documentclass[10pt,a4paper]{amsart}
\usepackage[margin=19mm]{geometry}
\usepackage{amsmath,amssymb,mathtools}
\usepackage[scr=rsfs,cal=euler]{mathalfa}
\usepackage{xfrac}
\usepackage[unicode,hidelinks,bookmarks=false]{hyperref}
\newcommand{\ie}{i.e.\ }
\newcommand{\cf}{cf.\ }
\newcommand{\resp}{respectively\ }
\newcommand{\Subsection}{\subsection}
\providecommand{\nobreakdash}{\nobreak\hbox{-}}
\setlength{\emergencystretch}{2em}
\multlinegap0pt
\usepackage{bm}
\usepackage{bbm}
\usepackage{dsfont}
\usepackage{xspace}
\usepackage{cancel}
\usepackage{mathtools}
\usepackage{amsthm}
\makeatletter
\@ifundefined{theorem}{\newtheorem{theorem}{Theorem}}{}
\@ifundefined{corollary}{\newtheorem{corollary}[theorem]{Corollary}}{}
\@ifundefined{lemma}{\newtheorem{lemma}[theorem]{Lemma}}{}
\makeatother
\newcommand{\gr}{\geqslant}
\newcommand{\ls}{\leqslant}
\newcommand{\vol}{\mathrm{vol}}
\title[On the maximal perimeter of isotropic log-concave probabilit]{On the maximal perimeter of isotropic log-concave probability measures}
\author{Silouanos Brazitikos, Apostolos Giannopoulos, Antonios Hmadi, Natalia Tziotziou}
\date{Source: arXiv:2602.03831v1 (CC BY 4.0)}
\begin{document}
\maketitle

\noindent\emph{Source and licence.} On the maximal perimeter of isotropic log-concave probability measures. Version arXiv:2602.03831v1 (CC BY 4.0), distributed under the Creative Commons Attribution 4.0 International licence, as recorded in the arXiv licence metadata for this exact version and shown on the abstract page https://arxiv.org/abs/2602.03831v1. Full rights notice: licenses/arxiv-2602.03831-CC-BY-4.0.txt.\par\smallskip

\noindent\textit{Editorial scope.} Counted unit: the Theorem whose source label is \texttt{th:perimeter} in the article (source0.tex lines 591--595, source label \texttt{th:perimeter}), delivered together with the article's own complete original proof of it (source0.tex lines 597--631, one \texttt{proof} environment of 35 lines). No mathematics is written, repaired or re-derived: statement and proof are the article's own text line for line. Required context retained uncounted: eq:surface-inradius (lines 254--256); lem:inradius (lines 393--396); cor:perimeter-subset-Rn (lines 576--579). Every definition, display and label of those lines is retained unchanged. Omitted from this package and not redistributed: 1--253, 257--392, 397--575, 580--590, 596--596, 632--1069. Nothing inside a retained excerpt is omitted or altered: every retained body is delivered exactly as the article's lines are written, so no omission receipt of any kind accompanies this package. Citations: the retained excerpts carry no citation command at all, so no bibliography entry of the article is transcribed and no reference is redistributed. Reference adaptation: none was needed and none was made; every reference inside the delivered unit resolves inside it, so no unresolved reference or citation is produced. Printed slips kept literal: no printed word, symbol, hypothesis or number of the counted statement or of its proof was altered. Layout and class: the article's own class is \texttt{article} with packages including srcltx, amsfonts, latexsym, amsmath, amssymb, amsthm, dsfont, fullpage, hyperref; the wrapper is the lane's pinned \texttt{amsart} editorial wrapper at 10pt on A4, which restates the theorem-like environments the retained text needs and adds the article's own macro definitions that the retained text needs (\texttt{\textbackslash gr}, \texttt{\textbackslash ls}, \texttt{\textbackslash vol}), with extra wrapper packages (bm, bbm, dsfont, xspace, cancel, mathtools). The delivered numbering is the unit's own. Assets: no retained span contains a figure environment, a graphics-inclusion command, a PDF-page inclusion, a table or a TikZ picture; no image file is copied into this package, the package declares no asset dependency and no image file is redistributed with it. Licence: version arXiv:2602.03831v1 of this article is distributed under the Creative Commons Attribution 4.0 International licence, as recorded in the arXiv licence metadata for this exact version and shown on the abstract page https://arxiv.org/abs/2602.03831v1. A full rights notice is delivered at licenses/arxiv-2602.03831-CC-BY-4.0.txt. Characters: every retained excerpt is pure ASCII, so the delivered engine, XeLaTeX with its default Latin Modern text font, reports no missing character.
\begin{equation}\label{eq:surface-inradius}
S(K)\ls\frac{n\vol_n(K)}{r(K)},
\end{equation}

\begin{lemma}\label{lem:inradius}
Let $\mu$ be an isotropic log-concave probability measure on $\mathbb{R}^n$, with $n\gr 3$. Then
$$r(R_{6n}(\mu))\gr 1/3.$$
\end{lemma}

\begin{corollary}\label{cor:perimeter-subset-Rn}Let $\mu$ be an isotropic log-concave probability measure on $\mathbb{R}^n$. For every  
convex body $A\subset\mathbb{R}^n$ with $A\subseteq R_{6n}(\mu)$, 
$$\mu^+(\partial A)\ls 14n^{3/2}.$$
\end{corollary}

\begin{theorem}\label{th:perimeter}Let $\mu$ be an isotropic log-concave probability measure on $\mathbb{R}^n$. For every  
convex body $A\subset\mathbb{R}^n$, 
$$\mu^+(\partial A)\ls Cn^{3/2},$$
where $C>0$ is an absolute constant.
\end{theorem}

\begin{proof}
Fix an arbitrary convex body $A\subset\mathbb{R}^n$. Note that if $A\subseteq R_{6n}(\mu)$ then the result follows from
Corollary~\ref{cor:perimeter-subset-Rn}. Otherwise, $A\setminus R_{6n}(\mu)\neq\varnothing$.  
 
We decompose $\partial A$ as follows:
$$\partial A=(\partial A\cap R_{6n}(\mu))\cup (\partial A\setminus R_{6n}(\mu)).$$
Note that $\partial A\cap R_{6n}(\mu)\subseteq \partial(A\cap R_{6n}(\mu))$, therefore 
\begin{equation}\label{eq:final-00}\mu^+(\partial A\cap R_{6n}(\mu))\ls \mu^+(\partial(A\cap R_{6n}(\mu)))\ls 14n^{3/2}\end{equation}
by Corollary~\ref{cor:perimeter-subset-Rn} applied to the convex body $A\cap R_{6n}(\mu)\subseteq R_{6n}(\mu)$.

Next, we estimate $\mu^+(\partial A\setminus R_{6n}(\mu))$. Note that
\begin{align}\label{eq:ell1-norm0}\|f\|_1 &=\int_0^{\|f\|_{\infty}}\vol_n(\{x\in\mathbb{R}^n:f(x)\gr t\})\,dt\\
\nonumber &=\int_0^{\infty}e^{-s}\|f\|_{\infty}\vol_n(\{x\in\mathbb{R}^n:f(x)\gr e^{-s}\|f\|_{\infty}\})\,ds\\
\nonumber &=\int_0^{\infty}e^{-s}\|f\|_{\infty}\vol_n(R_s(\mu))\,ds.\end{align}
We write
\begin{align*}
\mu^+(\partial A\setminus R_{6n}(\mu)) &=\int_{\partial A\setminus R_{6n}(\mu)}f(x)\,d\lambda(x)\\
&=\int_{\partial A\setminus R_{6n}(\mu)}\int_0^{\infty}\mathds{1}_{\{(x,t):f(x)\gr t\}}(x,t)\,dt\,d\lambda(x)\\
&=\int_0^{e^{-6n}\|f\|_{\infty}}\lambda (\partial A\cap \{x:e^{-6n}\|f\|_{\infty}\gr f(x)\gr t\})\,dt\\
&\ls\int_0^{e^{-6n}\|f\|_{\infty}}\lambda (\partial A\cap \{x:f(x)\gr t\})\,dt\\
&=\int_{6n}^{\infty}e^{-s}\|f\|_{\infty}\lambda (\partial A\cap R_s(\mu))\,ds,
\end{align*}
making the change of variables $t=e^{-s}\|f\|_{\infty}$. Since $\partial A\cap R_s(\mu)\subseteq \partial (A\cap R_s(\mu))$, and also $A\cap R_s(\mu)\subseteq R_s(\mu)$
and both sets are convex, using the monotonicity of 
Lebesgue surface area of convex bodies with respect to inclusion,
as well as \eqref{eq:surface-inradius} and Lemma~\ref{lem:inradius}, we get
$$\lambda (\partial A\cap R_s(\mu))\ls  \lambda (A\cap R_s(\mu))=S(A\cap R_s(\mu))\ls S(R_s(\mu))\ls \frac{n \vol_n(R_s(\mu))}{r(R_s(\mu))}
\ls 3n\vol_n(R_s(\mu)).$$
It follows that
\begin{equation}\label{eq:final-11}
\mu^+(\partial A\setminus R_{6n}(\mu))\ls 3n\int_{6n}^{\infty}e^{-s}\|f\|_{\infty}\vol_n(R_s(\mu))\,ds
\ls 3n\int_0^{\infty}e^{-s}\|f\|_{\infty}\vol_n(R_s(\mu))\,ds=3n\|f\|_1\end{equation}
using \eqref{eq:ell1-norm0} in the end. Combining \eqref{eq:final-11} with \eqref{eq:final-00} we obtain
$$\mu^+(\partial A)\ls 14n^{3/2}+3n\,\|f\|_1=(14+3/\sqrt{n})n^{3/2}.$$
This completes the proof of the theorem. \end{proof}

\end{document}
